\section{Algebraizing the equivariant deformation germ}
\label{sec:deformation}

We now construct the base \(\Bcore\) in
\cref{thm:main}.  The key point is that the invariant abstract tangent space
is a quotient of the invariant embedded tangent space.  It is not the
tangent space of the entire fixed Hilbert locus.

\subsection{Abstract and embedded unobstructedness}

Adjunction gives
\[
 K_{X_0}\simeq\mathcal O_{X_0}(-3),\qquad
 T_{X_0}\simeq\Omega^4_{X_0}(3).                   \tag{3.1}
\]
Akizuki--Kodaira--Nakano vanishing, in the form
\cite[Theorem 4.11]{Demailly1994}, gives
\[
 H^2(X_0,T_{X_0})
 =H^2(X_0,\Omega^4_{X_0}(3))=0,                   \tag{3.2}
\]
because \(4+2>5\).  The calculation may be made after base change to
\(\CC\); faithful base change then gives the asserted rational vanishing.
Thus the abstract Kuranishi germ is smooth.

Since \(X_0\) is a regular complete intersection,
\[
 N_{X_0/\PP^7}
 \simeq\mathcal O_{X_0}(2)\oplus\mathcal O_{X_0}(3).
\]
Kodaira vanishing yields
\[
 H^1(X_0,N_{X_0/\PP^7})=0.                         \tag{3.3}
\]
The Hilbert scheme is therefore smooth at \([X_0]\), with tangent
\[
 T_{\Hilb(\PP^7),[X_0]}=H^0(X_0,N_{X_0/\PP^7});
                                                               \tag{3.4}
\]
see \cite[Theorem 1.1(b),(c)]{Hartshorne2010}.

The restricted Euler sequence and Kodaira vanishing give
\[
 H^1(X_0,T_{\PP^7}|_{X_0})=0.
\]
Consequently the normal sequence induces a surjection
\[
 \kappa:H^0(N_{X_0/\PP^7})
 \twoheadrightarrow H^1(T_{X_0}).                  \tag{3.5}
\label{eq:3.5}
\]
For orientation, exact Hilbert-function calculations give
\[
 h^0(N)=35+111=146,\qquad
 h^0(T_{\PP^7}|_{X_0})=63.                         \tag{3.6}
\label{eq:3.6}
\]
These two numbers alone do not prove \(h^1(T_{X_0})=83\), because
\(H^0(T_{X_0})\) also occurs in the long exact sequence.

For an independent closure, the full linearized ideal-stabilizer equations
are
\[
 \delta_AQ_\rho=\nu Q_\rho,\qquad
 \delta_AC=\mu C+LQ_\rho.                          \tag{3.7}
\]
They form a coefficient matrix with \(156\) rows and \(74\) unknowns:
\(64\) entries of \(A\), the two scalars \(\nu,\mu\), and eight
coefficients of \(L\).  Exact row reduction has rank \(73\), and its kernel
is precisely
\[
 (A,\nu,\mu,L)=\lambda(I_8,2,3,0).                 \tag{3.8}
\]
Hence the projective infinitesimal stabilizer is zero,
\(H^0(T_{X_0})=0\), and \cref{eq:3.6} then gives \(h^1(T_{X_0})=83\).
This is a Lie-algebra calculation, not a classification of the full
\(\operatorname{PGL}_8\)-stabilizer.

\subsection{The fixed Hilbert point}

By \cref{lem:ambient-descent}, \(\Gscheme\) acts on a smooth Hilbert
neighborhood of \([X_0]\).  It is finite locally free of invertible order in
characteristic zero and therefore linearly reductive.  Romagny's fixed-point
smoothness theorem \cite[Theorems 1.2.1(2) and 4.3.6]{Romagny2022} implies
that the fixed-point germ
\[
 H^{\Gscheme}:=\Hilb(\PP^7)^{\Gscheme}
\]
is smooth at \([X_0]\).  Its tangent is
\[
 T_{H^{\Gscheme},[X_0]}=H^0(N)^{\Gscheme}.         \tag{3.9}
\]
Exactness of invariants turns \cref{eq:3.5} into
\[
 H^0(N)^{\Gscheme}\twoheadrightarrow
 H^1(T_{X_0})^{\Gscheme}.                          \tag{3.10}
\]

The exact Cayley computation in \cref{sec:cayley-local} gives
\[
 \dim_\Q H^1(T_{X_0})^{\Gscheme}=4.                \tag{3.11}
\]
Choose a rational complement
\[
 H^0(N)^{\Gscheme}=\ker(\kappa^{\Gscheme})\oplus L,
 \qquad \dim_\Q L=4.                               \tag{3.12}
\]
Because \(H^{\Gscheme}\) is smooth over \(\Q\) at a rational point, one
may choose an \'etale coordinate chart whose differential sends \(L\) to
the first four coordinate axes.  The inverse image of those axes is a
smooth locally closed rational four-germ \((\Bcore,0)\).  By construction,
\[
 T_{\Bcore,0}=L,\qquad
 \KS_0:T_{\Bcore,0}\xrightarrow{\sim}
 H^1(T_{X_0})^{\Gscheme}.                          \tag{3.13}
\label{eq:3.13}
\]

Restricting the Hilbert universal family gives
\[
 f:\mathcal X\longrightarrow\Bcore.
\]
The ambient action restricts to a fiberwise
\(\Gscheme_{\Bcore}\)-action.  Smoothness of the central fiber is open, so
we shrink the germ and obtain a smooth projective family.

\begin{remark}[Two algebraization firewalls]
The full fixed Hilbert germ has tangent \(H^0(N)^{\Gscheme}\) and retains
ambient-coordinate and equation-gauge directions; it is not asserted to
have dimension four.  Moreover, a tangent direction represented by a cubic
polynomial is not silently promoted to a literal family
\(C+\sum t_ip_i=Q_\rho=0\).  The actual family is the restricted Hilbert
universal family.
\end{remark}

After a splitting base change and an embedding into \(\CC\), the completed
family maps to the fixed Kuranishi germ.  Both germs are smooth of dimension
four and the derivative is \cref{eq:3.13}; hence the completed map is
formally \'etale.  Rim's equivariant versality theorem
\cite[p. 225]{Rim1980} supplies the compatible abstract equivariant
comparison, but it is not used to algebraize the family.
